\section{Better Unit Propagation}

So the textbook style unit propagation sucked. Let's change that. As discussed before, 
the biggest bottleneck/inefficiency in the actual implementation is the number of copies 
we need to make of our clauses and assignments on each decision level. Assigning one variable 
creates a new decision level, so 20 variables means 20 levels. For an unsatisfiable 
CNF equation, I estimate we need to make about 3 copies of our clauses and assignments:
one for the initial storage, another for the first failure, and a third for the second. 
Even if we ignore the inefficiencies in how we propagate (memory allocations, traversals, and 
adjustments our arrays would need to make when we delete things), this is a huge amount of 
memory to store. If you think that 20 variables and 91 clauses is a lot, modern SAT solvers 
can handle equations with \textbf{millions} of variables, so creating an algorithm that 
makes brute force look attractive is actually an achievement.
\footnote{They might revoke my degree for that actually}

So what's the strategy? Firstly, it would be nice if we didn't actually mutate our 
clauses at all. Ie, no clause deletion, and no variable deletion. Next, it would be 
nice if we didn't need hash sets in order to store our unit variables. Instead, it
would be nice if we had a data structure where we put in our decided variables, and then 
it slowly shrinks as we iterate through. Ie, we propagate one on variable at a time and 
never again. This sounds like we need a \textit{queue}.
\footnote{That wasn't shoehorned in at all!}

Lastly, instead of copying our assignment, and constantly reiterating through the 
\textit{entire} clause database, what if we just marked which variables we assigned,
and which clauses we've satisfied? That way when we need to undo, we just mark those 
variables as Unassigned, and those clauses as no longer satisfied. That way we save a lot 
of memory (and besides, copying is NOT O(1) as you might be led to believe by that deceptively
simply ``='' operator).

Let's change our interface:

\begin{lstlisting}
    public interface:
        SATSolver(string input)
        bool solveCNF()

    private interface:
        enum var
        {
            False
            True
            Unassigned
        }

        void parse(string equation)
        bool solve(var[] assignment, int level)
        bool evaluate(var[] assignment)
        (int[] vars, int[] clauses) propagate()

        int numVars
        int numClauses
        int numAssigned

        int[][] clauses
        bool[] satisfied

        optional var[] vars
\end{lstlisting}

We need to make some minimal changes to \textit{solveCNF()}:

\begin{algorithm}[H]
\caption{Main Solving Method}\label{alg:12}
\begin{algorithmic}[1]

\Procedure{solveCNF}{}

\State assignment[1..numVars]

\For{each $var$ in $assignment$}
    \State $var \gets Unassigned$
\EndFor

\State $\Call{Propagate}{assignment, 0}$

\State return $\Call{solve}{assignment, 1}$

\EndProcedure
\end{algorithmic}
\end{algorithm}

Since \textit{propagate()} now takes a decision level/variable index, we need to pass 
one to the method. Although no variable has an index of 0 (because it's reserved to mark the 
end of a clause), passing in $0$ just tells it ``Hey, there isn't a literal to propagate on,
so just find all clauses with only one variables and mark them as satisfied and `learn' those 
variables.'' And since this is at the root level, we don't need to undo anything if it turns out 
that the equation is unsatisfiable, since we're gathering things that HAVE to be a certain value 
if this equation is to be satisfiable.

Now let's get straight to \textit{solve()}:

\begin{algorithm}[H]
\caption{Recursive Solving}\label{alg:13}
\begin{algorithmic}[1]

\Procedure{solve}{assignment[1..n], level}

\If{level = numVars}
    \State return $\Call{evaluate}{assignment}$
\EndIf

\If{assignment[level] != Unassigned}
    \State $\Call{solve}{assignment, level + 1}$
\EndIf

\State $assignment[level] \gets true$
\State $(unitVars[1..n], satClause[1..m]) \gets \Call{propagate}{assignment, level}$

\If{$\Call{solve}{assignment, level + 1} = true$}
    \State return $true$
\EndIf

\For{each $var$ in $unitVars$}
    \State $idx \gets \Call{abs}{var}$
    \State $assignment[idx] \gets Unassigned$
\EndFor

\For{each $cIdx$ in $satClause$}
    \State $satisfied[cIdx] \gets false$
\EndFor

\State $assignment[level] \gets false$
\State $(unitVars[1..n], satClause[1..m]) \gets \Call{propagate}{assignment, level}$

\If{$\Call{solve}{assignment, level + 1} = true$}
    \State return $true$
\EndIf

\For{each $var$ in $unitVars$}
    \State $idx \gets \Call{abs}{var}$
    \State $assignment[idx] \gets Unassigned$
\EndFor

\For{each $cIdx$ in $satClause$}
    \State $satisfied[cIdx] \gets false$
\EndFor

\State $assignment[level] \gets Unassigned$
\State return $false$

\EndProcedure
\end{algorithmic}
\end{algorithm}

This method is practically the same, except now, we just check if we've already assigned 
the variable at the current decision level. If we have assigned a variable before we reach 
its decision level, then it must have been inferred/learned from the \textit{propagate()} 
method, so we don't need to call it again. Then, it's business as usual: assume the 
current variable is True, propagate, if it leads to a satisfiable assignment, return True, 
otherwise, we undo what we've assigned and marked as learned, then repeat with False. 

What's more interesting is our shiny new \textit{propagate()} method:

\begin{algorithm}[H]
\caption{Unit Propagation}\label{alg:14}
\begin{algorithmic}[1]

\Procedure{propagate}{assignment[1..n], level}

\State unit \Comment a queue of unit variables
\State unitVars[1..n]
\State satClauses[1..m]

\State $\Call{enqueue}{unit, level}$

\While{$Call{empty}{unit} = false$}
    \State $\Call{dequeue}{unit}$

    \For{$i \gets 1$ up to $numClauses$}
        \If{satisfied[$i$] = $true$}
            \texttt{Move onto the next clause}
        \EndIf

        \State $Clause[1..t] \gets Clauses[i]$
        \State $uNum \gets \Call{countUnassigned}{Clause}$
        \State $tNum \gets \Call{countTrue}{Clause}$

        \If{tNum > 0}
            \State satisfied[$i$] = $true$
            \State $\Call{append}{satClauses}{i}$
        \EndIf

        \If{tNum = 0 $\wedge$ uNum = 1}
            \State $\Call{append}{unitClauses, i}$
            \State $uVar \gets \Call{firstUnassigned}{Clause}$
            \State $\Call{append}{unitVars, uVar}$

            \State $idx \gets \Call{abs}{uVar}$
            \If{uVar > 0}
                \State $assignment[idx] \gets true$
            \Else
                \State $assignment[idx] \gets false$
            \EndIf

            \State $\Call{enque}{unit, uVar}$
            \State satisfied[$i$] = $true$
            \State $\Call{append}{satClauses, i}$
        \EndIf
    \EndFor
\EndWhile

\State return (unitVars, satClauses) 

\EndProcedure
\end{algorithmic}
\end{algorithm}

So, let's discuss what's going on. \textit{propagate()} now takes in the decision level, 
which is also the same thing as an index for a variable (conceptually speaking). It then 
adds that to a queue, then while the queue is not empty, it scans the clauses as before. 
First it makes sure we're not operating on a clause that's already been satisfied 
(because we can't tell anything from it), and when it finds a clause that has a true literal, 
it marks the clause as satisfied. Then when it finds a unit clause, it marks the clause as 
satisfied, keeps track of the new unit variable its found, and then adds that unit variable 
to the queue. Then once it's done, it returns the unit variables and the indices of the satisfied 
clauses.

You may have some questions about this approach. Firstly, how do we make sure we don't 
ever add the same variable twice? The answer is the same as before: because we assign a 
variable once it's found to be in a unit clause, it can't be found as Unassigned again, 
and therefore can't be marked as a unit variable again. Then, since we skip over already 
satisfied clauses, the amount of clauses we scan each iteration decreases monotonically. 

And going back to our \textit{solve()} code, even if we mark the same clause as satisfied 
or the same variable as unit twice, it doesn't matter because when we backtrack, we simply 
restore it back to its state of unsatisfied or Unassigned, respectively. What's important is 
that we don't mark the same clause or variable twice \textit{across decision levels} because 
then, and only then, would our backtracking code become a problem because we wouldn't be 
correctly restoring the state of the equation back to what it was before the decision level 
was entered. But since we can't mark the same variable or clause twice \textbf{until}
we backtrack, there's nothing to worry about.

Before we end this section, you may also be wondering: can we use the \textit{satisfied}
array in our \textit{evaluate()} method to skip clauses we know we've already satisfied?
And the answer is yes, and I would encourage you to do so in your implementations.

\subsection{Better Unit Propagation - Results}
So after implementing better unit propagation, this took \textbf{868ms} to solve all 1000 
satisfiable equations. The first time I implemented this version of unit propagation,
it actually took \textbf{2363ms}, so this is somehow over twice as fast. And because it's so 
fast, just like the first time I implemented this, we will be now testing our solver on 
unsatisfiable equations as well.

Our second testing set comes from the same database as before.
\footnote{\url{https://www.cs.ubc.ca/~hoos/SATLIB/benchm.html}, in case you needed a reminder.}
And our next testing set is called ``UUF50.218.1000." In other words, it's full of equations that 
have 50 variables, 218 clauses, and there are 1000 such equations in the set. In order to solve 
all 1000 equations, it took \textbf{140,411ms} (~2.34 minutes). That's also much faster than the first time 
I implemented this algorithm. The first time I did, it took 646,228ms (~10.77 minutes). Honestly,
with this much speedup, I don't know if I was just bad at implementing these algorithms the first time,
\footnote{Highly likely}
or if downgrading my operating system (and subsequently wiping my harddrive clean)
\footnote{I HATED MacOS Tahoe. Sequoia is where it's at.}
just has that much of an effect on execution speed.
\footnote{Also likely, depending on how much RAM and swap is being used to begin with.}

Why didn't we test our previous versions out on unsatisfiable equations? That's a good question.
For brute forcers, since there is no satisfiable assignment, it needs to check every single 
one to conclude that it's unsatisfiable, which would take a lot of time. Likewise with our 
recursive backtracking solver. And since our (my) implementation of textbook unit propagation 
was \textit{worse} than brute force,
\footnote{A difficult feat.}
obviously we wouldn't want to test it on something that will take more time than a satisfiable 
equation. Also because the asymptotic runtime is \textit{still} exponential, having a much 
larger equation would take a much, much larger amount of time, even if they were satisfiable.

\begin{FVerbatim}
Results so far:

On all 1000 equations in uf20-91:

Unit Propagation:          868ms (~0.87 seconds).
Recursive Backtracking:    1,540,613ms (~25.68 minutes).
Brute Force:               1,770,851ms (~29.51 minutes).
Short-Circuited (Basic):   49,874ms (~50 seconds).
Parallel:                  10,589ms (~10.5 seconds).

On all 1000 equations in UUF50.218.1000:

Unit Propagation: 140,041ms (~2.34 minutes).
\end{FVerbatim}